\documentclass[11pt,a4paper]{article}
\usepackage[utf8]{inputenc}
\usepackage[T1]{fontenc}
\usepackage[french]{babel}
\usepackage[margin=2.5cm]{geometry}
\usepackage{hyperref}
\usepackage{enumitem}
\usepackage{titlesec}
\usepackage{xcolor}

\definecolor{primary}{RGB}{37, 99, 235}
\titleformat{\section}{\color{primary}\Large\bfseries}{\thesection}{1em}{}[\vspace{-0.5em}\rule{\textwidth}{0.5pt}]
\titleformat{\subsection}{\large\bfseries}{\thesubsection}{1em}{}

\begin{document}

\begin{center}
    {\LARGE \bfseries Rapport de Projet -- Plateforme d'Annotation NLP}\\[0.5cm]
    {\large \textbf{Module :} Spring Boot \hfill \textbf{Année universitaire :} 2025/2026}\\[0.2cm]
    {\textit{Réalisé par : Amina Dourdi, Yousra Khallo}}\\[0.8cm]
\end{center}

\section*{Introduction}

Dans le cadre de notre formation, nous avons développé ce projet pour répondre à un besoin classique quand on fait de l'intelligence artificielle : la préparation des données. L'idée de base était de créer une plateforme où on peut importer facilement des textes, les distribuer à plusieurs personnes (qu'on appelle les annotateurs) pour qu'ils les lisent et les classent, et enfin utiliser ce travail pour entraîner un modèle de Machine Learning.

Pour concevoir quelque chose de simple et efficace, nous voulions une application avec deux côtés bien séparés. D'un côté, un espace administrateur pour importer les fichiers et gérer qui fait quoi. De l'autre, un espace annotateur hyper simple où l'utilisateur se connecte juste pour faire ses tâches sans se poser de questions. Techniquement, on a utilisé ce qu'on a vu en cours, principalement Spring Boot pour le backend, et on y a greffé du Python pour pouvoir lancer les entraînements de modèles. Ce rapport explique comment nous avons construit tout ça.

\section{Fonctionnalités développées}

\subsection{1. Partie Administrateur}

\textbf{Gestion des utilisateurs.} Quand l'admin crée un compte annotateur, un mot de passe aléatoire est généré côté serveur. Ce mot de passe part par e-mail grâce à \texttt{JavaMailSender} -- on a configuré un compte Gmail avec un mot de passe d'application pour ça. Au début, on envoyait juste le mot de passe et c'était tout, mais on s'est rendu compte que c'était un problème de sécurité : si quelqu'un intercepte le mail, il a un accès permanent. Du coup, on a rajouté un mécanisme de \textit{first login} : à la première connexion, une page s'affiche et force l'annotateur à choisir un nouveau mot de passe. Tant qu'il ne l'a pas fait, il ne peut rien faire d'autre.

\textbf{Import des datasets.} L'admin peut importer ses données sous forme de fichiers \textbf{CSV ou JSON}. Chaque fichier contient des paires de textes. Au moment de créer le dataset, il choisit aussi les classes d'annotation possibles -- dans notre cas de test, c'était \textit{entails}, \textit{neutral} et \textit{contradiction} pour du NLI, mais on peut mettre n'importe quelles classes.

\textbf{Affectation des tâches.} C'est probablement la partie qui nous a pris le plus de temps à réfléchir. On ne voulait pas que l'admin doive assigner les textes un par un. Nous avons donc écrit un algorithme qui prend la liste des annotateurs sélectionnés et distribue les couples de textes aléatoirement, en garantissant que chaque couple passe entre les mains d'\textbf{au moins 3 personnes} différentes. On avait besoin de ce minimum de 3 pour que le calcul du Kappa ait du sens ensuite. Si l'admin décide de retirer un annotateur en cours de route (parce qu'il est inactif, par exemple), le rééquilibrage se fait tout seul.

\textbf{Suivi de la progression.} On a ajouté une page dédiée où l'admin voit, pour chaque annotateur, combien de couples ont été traités et combien il en reste. C'est basique, mais c'est utile au quotidien pour savoir qui avance et qui traîne.

\textbf{Kappa de Fleiss et détection des spammeurs.} Pour vérifier que les annotateurs ne répondent pas n'importe quoi, on calcule le Kappa de Fleiss sur l'ensemble du dataset. Le score s'affiche dans une page Métriques. En parallèle, on a un mécanisme de détection des spammeurs : si un annotateur est trop souvent en désaccord avec les autres, il est signalé. L'admin peut alors décider de rejeter ses annotations. On a hésité entre un seuil fixe et un seuil configurable, et pour l'instant c'est un seuil fixe codé en dur.

\textbf{Entraînement ML.} L'admin peut lancer un entraînement directement depuis l'interface. En coulisses, Spring Boot appelle un script \texttt{train.py} qui charge le dataset annoté, vectorise les textes avec TF-IDF, et entraîne une Régression Logistique via \texttt{scikit-learn}. Les métriques (Accuracy, F1-Score) sont récupérées et affichées dans la page. On aurait aimé faire quelque chose de plus poussé avec des transformers, mais vu les délais, on est restés sur du classique.

\subsection{2. Partie Annotateur}

L'espace annotateur est volontairement plus simple. On voulait que la personne se connecte, tombe sur ses tâches, et commence à annoter sans avoir besoin d'explications.

Les couples de textes s'affichent un à la fois. L'annotateur lit les deux phrases, clique sur le bouton correspondant à la classe qu'il juge correcte, et passe au couple suivant. Quand il a tout terminé, la zone de saisie disparaît et un message lui confirme que c'est fini. On a fait ce choix pour éviter qu'il essaie de re-modifier ses réponses après coup.

Chaque annotateur a aussi accès à une page \textbf{profil} où il peut changer son nom, son prénom ou son mot de passe. Un tableau de bord lui montre ses statistiques : combien de tâches faites, combien en cours, son avancement en pourcentage.

Enfin, on a codé un \textbf{cron job} avec l'annotation \texttt{@Scheduled} de Spring. Tous les jours, cette tâche de fond vérifie les dates limites. Si un annotateur a une deadline le lendemain et qu'il n'a pas encore fini, il reçoit un mail de rappel automatiquement. On trouvait ça mieux qu'une simple notification dans l'interface parce que tout le monde ne se connecte pas forcément à la plateforme chaque jour.

\subsection{3. Sécurité et interface}

Pour la sécurité, on utilise \textbf{Spring Security} avec un système de rôles : \texttt{ADMIN} et \texttt{ANNOTATEUR}. Les mots de passe sont hachés en BCrypt. Côté frontend, on a stylisé les pages avec \textbf{Tailwind CSS}. On l'a chargé via CDN plutôt que de l'installer en local, c'était plus rapide à mettre en place.

\section{Améliorations possibles}

Si on devait continuer ce projet, il y a trois choses qu'on ferait en priorité. D'abord, passer sur \textbf{Docker} : actuellement, il faut installer Java, MySQL et Python séparément sur la machine, ce qui rend le déploiement un peu pénible. Ensuite, on remplacerait l'appel direct au script Python par une vraie \textbf{API REST} avec Flask ou FastAPI -- ça permettrait de brancher des modèles plus lourds (type BERT) sans bloquer l'application. Et enfin, on aimerait avoir plus de tests unitaires, surtout sur le calcul du Kappa et l'algorithme de répartition, parce qu'on a surtout testé ça manuellement pour l'instant.

\vspace{0.8cm}
\begin{center}
    \small\textit{Note : le \texttt{app.jar} fourni crée automatiquement les comptes \texttt{admin/admin} et \texttt{user1/user1} au lancement.}
\end{center}

\end{document}
